% ---------------------------------------------------------------------------
%  Shared preamble: page setup, palette, callout boxes, and the wiring toolkit.
%  Everything visual in the guide is drawn from these styles, so a wire is the
%  same colour in every figure it appears in.
% ---------------------------------------------------------------------------
\usepackage[T1]{fontenc}
\usepackage[utf8]{inputenc}
\usepackage{textcomp}
\usepackage{XCharter}
\usepackage[scaled=0.92]{FiraSans}
\usepackage[scaled=0.95]{zi4}
\usepackage{amssymb}
\usepackage[a4paper,top=22mm,bottom=20mm,left=22mm,right=22mm,headsep=7mm]{geometry}
\usepackage{graphicx}
\usepackage{booktabs}
\usepackage{array}
\usepackage{longtable}
\usepackage{multicol}
\usepackage{enumitem}
\usepackage{fancyhdr}
\usepackage{titlesec}
\usepackage{xcolor}
\usepackage{tikz}
\usepackage[most]{tcolorbox}
\usepackage[hidelinks]{hyperref}

\usetikzlibrary{calc,positioning,shapes.geometric,shapes.misc,arrows.meta,
                backgrounds,fit,decorations.pathreplacing,decorations.markings,
                patterns,shadows.blur}

% ── palette ────────────────────────────────────────────────────────────────
% The five signal colours are the ring's own colour vocabulary, taken from
% indicators.py, so a figure of the ring and a figure of a wire agree.
\definecolor{ink}{HTML}{1B1F24}
\definecolor{ink70}{HTML}{4A5057}
\definecolor{ink40}{HTML}{8A9199}
\definecolor{ink15}{HTML}{D8DDE2}
\definecolor{paper}{HTML}{FFFFFF}
\definecolor{shade}{HTML}{F4F6F8}

\definecolor{scred}{HTML}{C2185B}      % accent / raspberry
\definecolor{servo}{HTML}{E8710A}      % servo signal
\definecolor{buzzc}{HTML}{B26A00}      % buzzer
\definecolor{appc}{HTML}{00857A}       % the app, and the release it sends
\definecolor{btnc}{HTML}{0B7285}       % the take-dose button on the machine
\definecolor{dispc}{HTML}{9600FF}      % the screen, on either bus
\definecolor{rtcc}{HTML}{005AFF}       % rtc 3-wire (CE / I/O / SCLK)
\definecolor{p33}{HTML}{D24B2E}        % 3V3 rail
\definecolor{p5}{HTML}{B3001B}         % 5V rail
\definecolor{gndc}{HTML}{2B2B2B}       % ground
\definecolor{avoid}{HTML}{9AA3AB}      % deliberately not used

% The five states, as the LCD paints them. Same vocabulary the ring had.
\definecolor{ledgreen}{HTML}{00C83C}
\definecolor{ledblue}{HTML}{005AFF}
\definecolor{ledamber}{HTML}{FF8200}
\definecolor{ledcyan}{HTML}{00AABE}

\definecolor{okgreen}{HTML}{1E7A3C}
\definecolor{warnamber}{HTML}{9A6400}
\definecolor{dangerred}{HTML}{B3001B}

% ── shorthands for the units this guide keeps repeating ────────────────────
\newcommand{\ohm}{\,\textohm}
\newcommand{\uF}{\,\textmu F}
\newcommand{\dgr}{\textdegree}
\newcommand{\iic}{I\textsuperscript{2}C}
\newcommand{\code}[1]{\texttt{#1}}

% ── headings ───────────────────────────────────────────────────────────────
\titleformat{\section}
  {\sffamily\bfseries\Large\color{ink}}{\thesection}{0.6em}{}
  [\vspace{-0.6ex}{\color{ink15}\titlerule[1.1pt]}]
\titleformat{\subsection}
  {\sffamily\bfseries\large\color{ink}}{\thesubsection}{0.5em}{}
\titlespacing*{\section}{0pt}{2.2ex plus 1ex minus .2ex}{1.4ex plus .2ex}
\titlespacing*{\subsection}{0pt}{1.8ex plus 1ex minus .2ex}{0.9ex plus .2ex}

\setlength{\headheight}{14pt}
\pagestyle{fancy}
\fancyhf{}
\renewcommand{\headrulewidth}{0.4pt}
\renewcommand{\headrule}{\hbox to\headwidth{\color{ink15}\leaders\hrule height \headrulewidth\hfill}}
\fancyhead[L]{\sffamily\small\color{ink40}SnapCare pill dispenser \textendash{} assembly}
\fancyhead[R]{\sffamily\small\color{ink40}\leftmark}
\fancyfoot[C]{\sffamily\small\color{ink40}\thepage}

% ── code listings ──────────────────────────────────────────────────────────
\tcbuselibrary{listings,breakable,skins}
\lstdefinestyle{shellstyle}{
  basicstyle=\ttfamily\footnotesize\color{ink},
  breaklines=true, breakatwhitespace=false,
  columns=fullflexible, keepspaces=true,
  showstringspaces=false,
  commentstyle=\color{ink40}\itshape,
  morecomment=[l]{\#},
}
\newtcblisting{shell}{
  listing only, listing style=shellstyle,
  colback=shade, colframe=ink15, boxrule=0.6pt, arc=2pt,
  left=2mm, right=2mm, top=1.4mm, bottom=1.4mm,
  before skip=1.6ex, after skip=1.6ex,
}
\newtcblisting{pycode}{
  listing only,
  listing style=shellstyle,
  colback=shade, colframe=ink15, boxrule=0.6pt, arc=2pt,
  left=2mm, right=2mm, top=1.4mm, bottom=1.4mm,
  before skip=1.6ex, after skip=1.6ex,
}

% ── callout boxes ──────────────────────────────────────────────────────────
\newtcolorbox{provebox}[1][]{
  enhanced, breakable,
  colback=white, colframe=okgreen, boxrule=0pt,
  leftrule=2.4pt, arc=0pt, outer arc=0pt,
  left=3.5mm, right=3mm, top=2mm, bottom=2mm,
  borderline west={2.4pt}{0pt}{okgreen},
  title=\relax, #1
}
\newtcolorbox{warnbox}[1][]{
  enhanced, breakable,
  colback=warnamber!5, colframe=warnamber!45, boxrule=0.7pt, arc=2pt,
  left=3mm, right=3mm, top=2mm, bottom=2mm, #1
}
\newtcolorbox{dangerbox}[1][]{
  enhanced, breakable,
  colback=dangerred!5, colframe=dangerred!55, boxrule=0.9pt, arc=2pt,
  left=3mm, right=3mm, top=2mm, bottom=2mm, #1
}

\newcommand{\provetitle}{\textbf{\sffamily\color{okgreen}Prove it}\par\vspace{0.6ex}}
\newcommand{\warntitle}[1]{\textbf{\sffamily\color{warnamber}#1}\par\vspace{0.6ex}}
\newcommand{\dangertitle}[1]{\textbf{\sffamily\color{dangerred}#1}\par\vspace{0.6ex}}

% Caption styling
\usepackage[font={small,sf},labelfont={bf,color=ink70},textfont={color=ink70}]{caption}

% ── table helpers ──────────────────────────────────────────────────────────
\renewcommand{\arraystretch}{1.25}
\newcommand{\hdr}[1]{\textbf{\sffamily\small #1}}
\newcommand{\chk}{$\square$~}

% Colour swatch used in wire tables: a short length of the wire's own colour.
\newcommand{\swatch}[1]{\tikz[baseline=-0.5ex]{\draw[#1,line width=2.4pt,line cap=round] (0,0)--(0.5,0);}}

% ═══════════════════════════════════════════════════════════════════════════
%  THE 40-PIN HEADER
% ═══════════════════════════════════════════════════════════════════════════
% One macro, one data table. Every stage figure that shows a pin takes its
% colour from the same style list, so a pin is the same colour wherever it
% appears in the guide.
\tikzset{
  padbase/.style={rounded corners=1pt, draw=ink40, line width=0.5pt,
                  minimum width=6.4mm, minimum height=3.9mm, inner sep=0pt,
                  font=\sffamily\scriptsize},
  pfree/.style={padbase, fill=white, text=ink40},
  pavoid/.style={padbase, fill=avoid!18, text=ink40, draw=avoid},
  ppwr33/.style={padbase, fill=p33!18, text=p33!80!black, draw=p33!60},
  ppwr5/.style={padbase, fill=p5!14, text=p5!80!black, draw=p5!55},
  puse33/.style={padbase, fill=p33, text=white, draw=p33!70!black, line width=0.9pt},
  pusertc/.style={padbase, fill=rtcc, text=white, draw=rtcc!70!black, line width=0.9pt},
  puseservo/.style={padbase, fill=servo, text=white, draw=servo!70!black, line width=0.9pt},
  pusebuzz/.style={padbase, fill=buzzc, text=white, draw=buzzc!70!black, line width=0.9pt},
  pusebtn/.style={padbase, fill=btnc, text=white, draw=btnc!70!black, line width=0.9pt},
  pusedisp/.style={padbase, fill=dispc, text=white, draw=dispc!70!black, line width=0.9pt},
  pusegnd/.style={padbase, fill=gndc, text=white, draw=gndc, line width=0.9pt},
  plabel/.style={font=\sffamily\scriptsize, text=ink70, inner sep=1pt},
}

% \headerfigure{row spacing}  — draws the whole 40-pin header, top to bottom.
\newcommand{\drawheader}[1][0.62]{%
  \def\rowsep{#1}
  % board outline behind the pins
  \begin{scope}[on background layer]
    \fill[shade, rounded corners=2pt]
      (-4.9,0.62) rectangle (4.9,{-\rowsep*20-0.28});
  \end{scope}
  \node[font=\sffamily\bfseries\small, text=ink] at (0,0.36)
    {Raspberry Pi 3B \textendash{} 40-pin header};
  \node[font=\sffamily\scriptsize, text=ink70] at (0,0.06)
    {pin 1 is the corner nearest the microSD / power end};
  \foreach \r/\lb/\ls/\rb/\rs in {%
      1/{3V3}/{puse33}/{5V}/{ppwr5},
      2/{GPIO2 \textperiodcentered{} SDA}/{pusedisp}/{5V}/{ppwr5},
      3/{GPIO3 \textperiodcentered{} SCL}/{pusedisp}/{GND}/{pusegnd},
      4/{GPIO4 \textperiodcentered{} CTP INT}/{pusedisp}/{GPIO14 TXD}/{pavoid},
      5/{GND}/{pusegnd}/{GPIO15 RXD}/{pavoid},
      6/{GPIO17}/{puseservo}/{GPIO18 \textperiodcentered{} LCD RST}/{pusedisp},
      7/{GPIO27}/{puseservo}/{GND}/{pusegnd},
      8/{GPIO22}/{puseservo}/{GPIO23}/{pusebuzz},
      9/{3V3}/{ppwr33}/{GPIO24}/{pusebtn},
      10/{GPIO10 \textperiodcentered{} MOSI}/{pusedisp}/{GND}/{pusegnd},
      11/{GPIO9 MISO}/{pavoid}/{GPIO25 \textperiodcentered{} DC}/{pusedisp},
      12/{GPIO11 \textperiodcentered{} SCK}/{pusedisp}/{GPIO8 \textperiodcentered{} CS}/{pusedisp},
      13/{GND}/{pusegnd}/{GPIO7 CE1}/{pavoid},
      14/{GPIO0 ID}/{pavoid}/{GPIO1 ID}/{pavoid},
      15/{GPIO5}/{puseservo}/{GND}/{pfree},
      16/{GPIO6}/{puseservo}/{GPIO12 \textperiodcentered{} CTP RST}/{pusedisp},
      17/{GPIO13}/{puseservo}/{GND}/{pfree},
      18/{GPIO19}/{puseservo}/{GPIO16 \textperiodcentered{} SCLK}/{pusertc},
      19/{GPIO26}/{puseservo}/{GPIO20 \textperiodcentered{} I/O}/{pusertc},
      20/{GND}/{pusegnd}/{GPIO21 \textperiodcentered{} CE}/{pusertc}%
    }{%
    \pgfmathtruncatemacro{\lp}{2*\r-1}%
    \pgfmathtruncatemacro{\rp}{2*\r}%
    \pgfmathsetmacro{\yy}{-\rowsep*\r}%
    \node[\ls] (L\lp) at (-0.62,\yy) {\lp};
    \node[\rs] (R\rp) at (0.62,\yy) {\rp};
    \node[plabel, anchor=east] at (-1.05,\yy) {\lb};
    \node[plabel, anchor=west] at (1.05,\yy) {\rb};
  }%
}

% A compact legend for the header figure.
\newcommand{\headerlegend}{%
  \begin{tikzpicture}[every node/.style={font=\sffamily\scriptsize,text=ink70}]
    \def\sw##1##2##3{\node[##1,minimum width=4.6mm,minimum height=3.2mm] at (##2,0) {};
                     \node[anchor=west] at (##2+0.28,0) {##3};}
    \sw{puseservo}{0}{servo signal}
    \sw{pusedisp}{2.9}{screen}
    \sw{pusebuzz}{5.6}{buzzer}
    \sw{pusebtn}{7.4}{button}
    \begin{scope}[yshift=-0.52cm]
      \sw{pusertc}{0}{RTC 3-wire}
      \sw{puse33}{2.9}{3V3}
      \sw{pusegnd}{5.6}{ground}
      \sw{pavoid}{7.4}{do not use}
    \end{scope}
  \end{tikzpicture}%
}

% ═══════════════════════════════════════════════════════════════════════════
%  WIRING-FIGURE TOOLKIT  —  used by every stage
% ═══════════════════════════════════════════════════════════════════════════
\tikzset{
  modbox/.style={draw=ink70, line width=0.8pt, fill=white, rounded corners=2pt,
                 minimum width=30mm, minimum height=14mm, align=center,
                 font=\sffamily\small},
  modtitle/.style={font=\sffamily\bfseries\small, text=ink},
  term/.style={circle, draw=ink70, fill=white, line width=0.5pt,
               inner sep=0.6pt, minimum size=2.6mm, font=\sffamily\tiny},
  wire/.style={line width=1.5pt, line cap=round, rounded corners=2pt},
  compbox/.style={draw=ink40, line width=0.6pt, fill=shade, rounded corners=1.5pt,
                  font=\sffamily\scriptsize, inner sep=2.4pt, align=center},
}

\renewcommand{\sectionmark}[1]{\markboth{#1}{#1}}

% A single header pad in a stage figure. Named H<pin>, so wires can say (H11).
% #1 x  #2 y  #3 physical pin  #4 pad style  #5 label, placed to the right
\newcommand{\hpin}[5]{%
  \node[#4] (H#3) at (#1,#2) {#3};
  \node[plabel, anchor=west] at ({#1+0.42},#2) {#5};
}

% A frame with a caption behind a strip of pads.
\newcommand{\hframe}[4]{% #1 x #2 ytop #3 ybottom #4 caption
  \begin{scope}[on background layer]
    \fill[shade, rounded corners=2pt]
      ({#1-0.55},{#2+0.34}) rectangle ({#1+2.5},{#3-0.28});
  \end{scope}
  \node[font=\sffamily\scriptsize\bfseries, text=ink40, anchor=west]
       at ({#1-0.5},{#2+0.62}) {#4};
}


% resistor drawn inline on a wire between two coordinates
\newcommand{\inlineres}[4]{% #1 colour #2 from #3 to #4 label
  \coordinate (rA) at ($(#2)!0.38!(#3)$);
  \coordinate (rB) at ($(#2)!0.62!(#3)$);
  \draw[wire,#1] (#2) -- (rA);
  \draw[wire,#1] (rB) -- (#3);
  \draw[ink70,line width=0.7pt,fill=white] ($(rA)!0.5!(rB)$) ellipse (0.34 and 0.13);
  \node[font=\sffamily\tiny,text=ink] at ($(rA)!0.5!(rB)$) {#4};
}
